\section{Introduction}\label{sec:intro}

The loss ratio, claims divided by premium, is among the most widely used quantities in insurance
underwriting, portfolio management and reinsurance pricing. Ex post it is incurred claims over earned
premium; ex ante it is the ratio of a loss measure $\rho(X)$, such as an expected or a risk-adjusted loss, to a
premium $\Pi(X)$ of the same position $X$. When the premium is a non-linear functional of the loss, or the
numerator is a risk measure, it is not clear which functionals of a loss should be called a loss ratio, or
which properties a loss ratio should have when books are pooled.

The question arises in risk reporting. Before the claims are known, a loss ratio is usually reported through
its expected value, expected claims over premium. When the risk of a book is concentrated in the tail of its
loss distribution, the expected loss ratio can be misleading: two books with the same expected loss ratio may
differ greatly in their exposure to large losses. A natural remedy is to replace the mean of the claims by a tail
measure, and to report, for instance, the value at risk or the tail value at risk (expected shortfall) of the claims divided by the
premium \citep{adeh1999}. The loss ratio then becomes a model-based, or mark-to-model, quantity: a risk measure
of the loss, evaluated under a model of its distribution, divided by a premium principle. The ex-post accounting
ratio is the special case in which the model is the realized outcome. Once such ratios are admitted, the
question is which properties are necessary and sufficient for a functional of the loss to be a loss ratio of
this kind; this is the question addressed here. The tail value at risk and the value at risk are distortion
functionals, so tail loss ratios over a distortion premium are covered by the results below
(see the ratio $\mathrm{ES}_{0.9}/\E$ in Example~\ref{ex:qc}).

We work with the \emph{extended loss ratio} $\LR=\rho/\Pi$, which allows negative values and zero premiums,
defined abstractly by \emph{association}: a family of risk measures $\{\rho_\beta\}$ is
associated with $\LR$ if $\LR(X)=\beta\iff\rho_\beta(X)=0$; the canonical family is
$\rho_\beta=(\rho-\beta\Pi)/(1-\beta)$, $\beta\ne1$. The present paper characterizes, by axioms stated on $\LR$
alone, the loss ratios that are ratios of two Choquet-type functionals.

Three features distinguish the problem from the classical characterizations of premium principles and risk
measures \citep{yaari1987,wyp1997,schmeidler1986,www2020}. First, the object is a ratio, while those results
concern a single functional. Second, the ratio does not determine its numerator and denominator: on positive
losses, $(\rho,\Pi)$ and $(\psi\rho,\psi\Pi)$ give the same ratio for every positive functional $\psi$. Third, the
ratio is not monotone: we show that a loss ratio of two distortion functionals is monotone only if it is
identically $1$ (Proposition~\ref{prop:F2}), so the theory of acceptability indices, which rests on monotonicity
\citep{chernymadan2009}, does not apply.

Our approach is to recover the premium and the risk measure from the loss ratio. For a book $X$, consider the
\emph{cash line} $c\mapsto\LR(X+c)$. A premium is the sure amount at which a risk changes hands; on the cash line,
a zero premium is a point at which the loss ratio is not finite. We define the \emph{implied premium} $\Pi(X)$ as
minus the unique such point, and the \emph{implied risk measure} by the accounting identity
$\rho(X)=\Pi(X)\LR(X)$, extended by cash shifts to books with zero premium. At such books the implied risk
measure is the residue of the cash line at its singularity (Proposition~\ref{prop:readings}). For a Bernoulli loss
$\1_A$ with $P(A)=\tfrac14$, the expected loss and the square-root distortion premium, the cash line is
$(\tfrac14+c)/(\tfrac12+c)$; it is singular at $c=-\tfrac12$, so $\Pi(\1_A)=\tfrac12$ and
$\rho(\1_A)=\tfrac12\LR(\1_A)=\tfrac14$ (Figure~\ref{fig:three}).

With this reading, Theorems~\ref{thm:A} and~\ref{thm:B} state the following. A map $\LR$ from bounded losses to
$\R\cup\{\infty,\star\}$ is of the form $I_h/I_g$, with $I_h,I_g$ signed Choquet integrals of distortions $h,g$
of bounded variation, if and only if it satisfies the following axioms, stated in Section~\ref{sec:axioms}:
\axP{} every cash line has exactly one point where $\LR$ is not finite; \axN{} $\LR(c)=1$ for sure losses
$c\ne0$; \axCM{} the loss ratio of a pool of comonotone books is a combination of their loss ratios with their relative costs, $\Pi(X)/(\Pi(X)+\Pi(Y))$, as weights; \axLI{} $\LR$ is
law invariant; and \axRC{} the implied premium and risk measure are uniformly continuous. If \axRC{} is replaced
by \axRM, monotonicity of the implied premium and risk measure, $h$ and $g$ are distortion functions and no continuity
assumption is needed. In both cases $(h,g)$ is unique, $\Pi=I_g$ and $\rho=I_h$. The unit-slope Möbius form of the
cash lines, the normalization $\Pi(c)=c$, comonotonic additivity of the implied premium, scale invariance and
the association are consequences of the axioms.

Given \axP{} and \axN, axiom \axCM{} is equivalent to comonotonic additivity of the implied premium and risk measure,
the axiom of \citet{yaari1987} and \citet{wyp1997}; once this equivalence is established, the representation step
follows from the theory of comonotonic additive functionals \citep{schmeidler1986,msm1994,cmmm2012,www2020}. The
contributions specific to loss ratios are: the construction of the implied premium and risk measure; the proof
that the pooling identity yields additivity of the implied pair on all comonotone pairs, including books that
break even, on which a premium given exogenously would not be identified (Lemma~\ref{lem:comon} and
Remark~\ref{rem:primitive}); unconditional uniqueness of $(h,g)$; the impossibility of a monotone loss ratio and
its consequences for the associated family (Corollary~\ref{cor:assoc}); a version for nonnegative losses
(Theorem~\ref{thm:C}); and counterexamples showing that none of the axioms
\axN, \axCM, \axLI, \axRC, \axRM{} can be dropped.

The paper has a second aim: it is a brief case study of the use of AI in mathematical research in actuarial
science. The work was done with Claude (Anthropic), used as a main assistant and to coordinate panels of AI
agents on designated parts of the project: a roundtable of six agents and a larger panel covering ten literature
strands surveyed the literature, proposed candidate statements, ranked them by judge agents, drafted proofs and
submitted them to repeated AI-assisted critical review; the main assistant constructed examples and
counterexamples and drafted and revised the manuscript. ChatGPT (OpenAI) was used for an informal review of the
proofs, and Gemini (Google) for summaries of key cited works. The problem was chosen by the author, from the
risk-reporting question above, and in part as a deliberate test: its solution requires machinery, such as
signed Choquet integrals, that is not routine, which makes it a demanding case for AI assistance. The research questions and all decisions on content
were the author's, and the formulation through the implied premium and the implied risk measure arose from the
author's objections to earlier AI-drafted versions, in particular that an assumed cash-line axiom supplied $\rho$
and $\Pi$ by fiat. The AI-generated material was treated as proposals: the reviews found errors, including a
wrong sufficient condition for quasi-convexity and an incorrect uniqueness claim, which were corrected (Table~\ref{tab:errors}); the
numbers in the examples were recomputed; the statements taken from three of the main sources were checked against
those papers; and the steps specific to loss ratios were formalized in the proof assistant Lean and checked by its
kernel, with the classical representation step taken as a hypothesis. Section~\ref{sec:collab} describes this
process, the errors and guardrails, and what has not been verified, and illustrates a standard for reporting such
use in a mathematical paper.

The paper is organized as follows. Section~\ref{sec:lit} reviews the related literature, and
Section~\ref{sec:setting} introduces the extended loss ratio. Section~\ref{sec:axioms} defines the implied premium
and risk measure and states the axioms. Section~\ref{sec:results} contains the main theorems and examples, and
Section~\ref{sec:triviality} discusses the role of the axioms. Proofs are in Section~\ref{sec:proofs}.
Section~\ref{sec:actuarial} treats nonnegative losses, Section~\ref{sec:necessity} gives the counterexamples, and
Section~\ref{sec:consequences} derives consequences. Section~\ref{sec:discussion} lists open problems, and
Section~\ref{sec:collab} describes the use of AI, including the formal check in Lean
(Section~\ref{sec:collab-lean}).
